\chapter{通往主动推断的高路}

\begin{quote}
能够模拟未来的生存机器，比那些只能依赖公开的试错来学习的生存机器领先一步。公开试验的问题在于它需要时间和能量。公开犯错的问题在于它往往是致命的。模拟既更安全，也更迅速。

\hfill --- 理查德\textperiodcentered 道金斯
\end{quote}

\section{引言}

在第 2 章中，我们把自由能的引入说明为一种执行近似贝叶斯推断的方法（也就是通往主动推断的“低路”）。本章则从另一种视角引入自由能，即“高路”：它倒转了前面的论证顺序，从统计物理学的第一性原理以及生物体必须维持自身存在这一核心要求出发，也就是说，生物体必须避免令人惊讶的状态，然后再把自由能最小化引入为这一问题在计算上可处理的解法。本章揭示了变分自由能最小化与近似贝叶斯推断中模型证据最大化（或称“自证”）之间的形式等价性，从而展示了自由能与关于适应系统的贝叶斯视角之间的联系。最后，本章讨论主动推断如何提供一种新的第一性原理视角，用来理解（最优）行为。

主动推断是一种理论，说明生物有机体如何通过知觉与行动最小化惊讶，或者最小化惊讶的一个可处理替代量，即变分自由能，从而维持自身存在。通过从第一性原理出发，它推进了一种基于信念的新方案，用以理解行为与认知，而这一方案带来了大量经验层面的含义。

通往主动推断的高路始于这样一个前提：为了生存，任何生物体都必须把自身维持在一组适当的偏好状态之中，同时避开环境中的其他非偏好状态。这些偏好状态首先由生态位特定的进化适应所界定。不过，正如我们稍后将看到的，在高级生物体中，这些状态也可以扩展到习得的认知目标。例如，为了生存，一条鱼必须停留在一个舒适区内，而这个舒适区只对应宇宙所有可能状态中的一个很小子集：它必须待在水里。类似地，人类必须确保自己的内部状态（例如体温、心率之类的生理变量）始终保持在可接受范围内，否则他们就会死亡（更准确地说，会变成别的东西，比如一具尸体）。这个可接受的范围或舒适区，以约定方式界定了某个对象要成为“它自身”所必须处于的特征状态。

生物体通过在多个层次上主动控制自身状态来解决这一根本的生物学问题，例如调节体温。这些层次从出汗这样的自动调节机制（生理学），延伸到购买并饮用一杯饮料这样的认知机制（心理学），再到配置空调系统这样的文化实践（社会科学）。

从更形式化的角度看，主动推断把生存这一生物学问题，或者说生存的解释，表述为惊讶最小化。这一表述建立在信息论中关于“令人惊讶状态”的技术性定义之上，本质上，令人惊讶的状态就是那些落在生物体舒适区之外的状态。随后，它提出以自由能最小化作为一种实践上可行、并且具有生物学基础的方式，使生物体或适应系统最小化其感官遭遇的惊讶。

\section{马尔可夫毯}

任何适应系统的一个重要前提，是它必须与环境享有某种分离与自主性，否则它就会简单地耗散、溶解，并因此屈从于环境动力学。没有这种分离，就不会有需要最小化的惊讶；必须有某个会感到惊讶的东西，以及某个使其感到惊讶的东西。换言之，至少存在两个东西，即系统与环境，而且二者可以彼此区分。表达系统与其余环境之间分离的一种形式化方式，是马尔可夫毯这一统计构造（Pearl 1988）；见框 3.1。

\paragraph{框 3.1 马尔可夫毯}
马尔可夫毯是本书中一个重要且反复出现的概念（Friston 2019a, Kirchhoff et al. 2018, Palacios et al. 2020）。从技术上说，毯状态 $b$ 的定义如下：

\[
\mu \perp x \mid b \iff p(\mu, x \mid b) = p(\mu \mid b)p(x \mid b)
\]

这句话用两种不同但等价的方式说明：如果已知 $b$，那么变量 $\mu$ 与变量 $x$ 条件独立。换句话说，如果我们知道 $b$，那么再知道 $x$，也不会给我们带来关于 $\mu$ 的额外信息。一个常见例子是马尔可夫链：过去导致现在，现在导致未来。在这种情形下，过去只能通过现在影响未来。这意味着，在已知现在的前提下，再去了解过去，并不会给我们带来关于未来的额外信息。

如果我们已知某个系统中的条件依赖关系，那么要识别其中的一个马尔可夫毯，可以遵循一个简单规则：某个给定变量的毯，包含它的父节点（即它所依赖的变量）、它的子节点（即依赖于它的变量），以及在某些设定下，它那些子节点的其他父节点。

简而言之，马尔可夫毯是一组变量，它中介了系统与其环境之间全部的统计交互。图 3.1 展示了马尔可夫毯在动态场景中的一种解释。在这里，条件独立关系又补充了动力学约束，因此流不会依赖于毯另一侧的状态。

图 3.1 中的马尔可夫毯，把适应系统内部的状态（即大脑活动）与环境的外部状态区分开来。此外，它还识别出另外两类状态，分别标记为感觉状态与主动状态，这两类状态共同构成了在统计上把内部状态与外部状态分开的“毯”。所谓统计分离，意味着如果我们知道主动状态和感觉状态，那么外部状态就不会再为内部状态提供额外信息，反之亦然。在动力学语境中，这通常被解释为：内部状态不能直接改变外部状态，但可以通过改变主动状态而间接地做到这一点；同样，外部状态不能直接改变内部状态，但可以通过改变感觉状态而间接地做到这一点。

这其实是经典行动--知觉循环的重述：适应系统与环境分别只能通过行动与观测发生交互。这种重新表述有两个主要好处。

\paragraph{图 3.1}
一个动态马尔可夫毯，它把一个适应系统（此处是大脑）与环境分隔开来。每组状态的动力学都由一个确定性流决定，这个流由函数 $f$ 指定，给出平均变化率，并叠加额外的随机波动 $\omega$。箭头表示各变量对其他变量变化率的影响方向（从技术上说，就是相关雅可比矩阵中的非零元素）。这只是一个例子；马尔可夫毯既可以用来把整个有机体与环境分开，也可以彼此嵌套多个马尔可夫毯。例如，大脑、有机体、二人关系以及共同体，都可以被理解为不同层级、彼此嵌套的马尔可夫毯（形式化处理见 Friston 2019a；Parr, Da Costa, and Friston 2020）。容易让人困惑的是，不同领域对变量采用不同记号；有时感觉状态记作 $s$，外部状态记作 $\eta$，主动状态记作 $a$。这里我们选择当前这组变量，是为了与本书其他章节保持一致。

第一，它形式化地表达了这样一个事实：适应系统的内部状态对于环境动力学具有自主性，因此能够抵抗环境影响。第二，它为适应系统如何最小化惊讶提供了结构支架：它凸显出系统能够接触到的内部状态、感觉状态和主动状态。具体而言，惊讶是相对于感觉状态来定义的，而内部状态与主动状态的动力学，则是最小化感觉状态惊讶的手段。

这里需要注意的关键点是，适应系统的内部状态与外部状态之间存在一种形式关系。这源于马尔可夫毯两侧的一种对称性，因为二者都会影响毯状态，同时也都会被毯状态影响。由此带来的一个结果是，我们可以在给定毯状态的条件下，为内部状态和外部状态分别构造条件概率分布。由于这两个分布都条件化于同一组毯状态，我们就可以把成对的期望内部状态与期望外部状态联系起来。换言之，平均而言，内部状态与外部状态获得了一种（广义的）同步，就像我们在一根木梁两端各挂一个钟摆时所预期的那样。随着时间推移，当它们逐渐同步时，每个钟摆都会通过木梁的间接影响而对另一个钟摆具有预测性（Huygens 1673）。图 3.2 为这种关系提供了图示直觉。这意味着，如果我们能够写下在给定马尔可夫毯条件下外部状态与内部状态各自独立的分布，那么这两类状态就会经由这层毯而彼此提供信息。

这种同步使内部状态看起来像是在表征（或建模）外部状态，这又把我们带回第 2 章引入的惊讶最小化思想。因为惊讶依赖于一个关于感觉数据如何生成的内部模型。回顾一下，对感觉观测之惊讶（负对数概率）的最小化，与对该模型证据（边缘似然）的最大化是同一回事，而模型证据就是在该模型下感觉观测出现的概率。对惊讶最小化的这一观念，可以从两个等价视角来理解：贝叶斯视角与自由能视角。下面我们将讨论这两种视角。

\section{惊讶最小化与自证}

从贝叶斯视角看，一个拥有马尔可夫毯的主体似乎在对外部环境进行建模，因为其内部状态（平均而言）对应于对系统外部状态的一种概率表征，也就是一种近似后验信念（图 3.2）。内部状态的动力学对应于对外部状态进行的一种（近似）贝叶斯推断，因为内部状态的运动会改变相应的概率分布，而这种变化又得益于一个隐含的生成模型，该模型说明感觉，或者用马尔可夫毯术语来说“感觉状态”，是如何产生的。如果我们恢复“主体由内部状态和毯状态构成”这一观念，那么我们就可以谈论主体的生成模型。

重要的是，主体的生成模型不能只是简单复制外部动力学，否则主体就只会跟随外部的耗散动力学。相反，这个模型还必须规定主体得以存在的偏好条件，或者主体为了维持自身存在而必须访问的状态区域，或者说，以占据特征状态的方式来满足其存在标准。这些偏好状态（或观测）可以被指定为模型的先验，这意味着模型隐含地假定：如果主体满足了自身存在的标准，那么它偏好的（先验的）感觉更有可能出现，也就是更不令人惊讶。这意味着模型带有一种隐性的乐观偏差。这种乐观偏差是必需的，因为只有这样，主体才不会止于复制外部动力学，而是能够规定那些支撑其偏好状态或特征状态的主动状态。

在这种表述下，可以把最优行为（相对于先验偏好而言）理解为通过知觉与行动实现模型证据最大化。事实上，模型证据概括了生成模型对感觉的拟合或解释程度。良好的拟合表明，该模型成功说明了自己的感觉输入（这是推断的描述性一面）；与此同时，它也实现了自己的偏好感觉，因为这些感觉更不令人惊讶（这是推断的规定性一面）。这种良好拟合保证了惊讶最小化，因为最大化模型证据 $P(y)$ 在数学上等价于最小化惊讶：$\Im(y) = -\ln P(y)$。

对上述论证的一种更简洁重述，是说任何适应系统都在进行“自证”（self-evidencing）（Hohwy 2016）。这里的自证，是指通过行动去获取与内部模型一致的感觉数据，也就是为该模型提供证据，从而最大化模型证据。

\subsection{作为哈密顿最小作用量原理的惊讶最小化}

在前面的几节中，我们断言惊讶必须被最小化，但还没有详细说明原因。尽管自证背后的物理学细节超出了本书的范围（详见 Friston 2019b），这里我们仍对这些原理作一个简要概述。这些原理建立在这样一个思想之上：带有马尔可夫毯的生物体能够在时间中持续存在，并抵抗环境波动的离散效应。马尔可夫毯的持续存在意味着，毯状态的分布在时间中保持恒定。简单地说，这意味着每当感觉状态（或主动状态）偏离该分布下高概率区域时，这种偏离都必须由状态的平均流来纠正，而平均流正是图 3.1 中流的确定性部分。如果用物理学家的方式来表述，那么处于稳态的随机系统会进行这样的动力学：它们平均而言沿着一个能量函数（或哈密顿量）下降，而这个能量函数可以被解释为负对数证据或惊讶。这就像一个球从山顶的高重力势能处滚下，最后滚入低能量的盆地。见图 3.3。

对于图 3.3 左侧所示的系统，每当一次波动导致系统进入一个概率更低的状态，这种偏移都会通过沿概率梯度上升的运动得到纠正，因此该系统会在更大比例的时间里占据概率密集的区域。这里的关键洞见在于：这个系统通过最小化惊讶（平均而言）把感觉状态维持在一个狭窄范围内；相比之下，右侧那个系统的惊讶则会无限增长。

惊讶最小化使得生物体能够（暂时地）抵抗热力学第二定律。第二定律指出，熵，也就是系统状态的离散程度，总是在增长。这是因为，平均而言，熵就是惊讶的长期平均，而平均来说，最大化观测对数概率等价于最小化香农熵：

\begin{equation}
H[P(y)] = E_{P(y)}[\Im(y)] = -E_{P(y)}[\ln P(y)]
\tag{3.1}
\end{equation}

确保只有一小部分感觉状态以高概率被占据，就等价于维持某种特定熵。这正是自组织系统的一个定义性特征，也是控制论理论早已认识到的事实。

从生理学视角看，惊讶最小化形式化了稳态（homeostasis）这一思想。当某个传感值偏离其最优范围时，负反馈机制就会启动，并逆转这些偏离。从控制论视角看，我们可以把最优行为理解为相对于某种目标稳态概率密度而言的行为。换言之，如果我们定义一个偏好结果的分布，那么最优行为就会推动系统向该分布演化，并维持在该分布附近。

正如第 2 章所见，自由能是惊讶的一个上界，这意味着面对随机波动时，可以通过最小化自由能来获得最优行为。回忆一下，自由能与惊讶之间的差异，是精确后验概率（即在给定毯状态下外部状态的分布）与近似后验概率（即在给定平均内部状态下外部状态的分布）之间的散度。因此，内部状态的运动可以被理解为在最小化这一散度，而这又使主动状态平均而言能够最小化伴随感觉状态出现的惊讶。换句话说，由自由能最小化导出的最优行为，是最不令人惊讶的行为，并且它沿着一条从当前状态通向目标状态的最小作用量路径前进，也就是说，这是把哈密顿最小作用量原理应用到了行为之上。

图 3.3 给出了一个非常简单的例子：一个具备随机吸引子的系统。这类似于一个恒温器，它在控制论语言中只有单一设定点，既不能学习，也不能规划。主动推断试图使用同一套解释装置，去覆盖复杂得多且适应性强得多的系统。在这里，最简单系统与更复杂系统之间的区别，可以归结为它们吸引子形状的不同，从固定点到越来越复杂、越来越游走的动力学。由此可见，生物体可以被理解为不断在过度稳定与过度离散之间寻求折中，而主动推断的目标，就是解释这种折中是如何实现的。

\section{推断、认知与随机动力学之间的关系}

物理学家 E. T. Jaynes 曾著名地指出，推断、信息论与统计物理学，只是对同一事物的不同视角（Jaynes 1957）。前面几节中，我们讨论了贝叶斯视角与统计物理学视角如何为惊讶最小化和最优行为提供两种等价的理解方式，这实际上等于在 Jaynes 的三元组中又加入了一种认知维度。各种思想流派之间的这种等价性很有吸引力，但对不熟悉相应形式体系的人来说也可能令人困惑，因为许多不同术语指向的其实是相同的量。为了帮助澄清这一点，本节将详细说明贝叶斯视角与统计物理学视角之间的主要等价关系，以及它们在认知上的解释；概览见表 3.1 和框 3.2。

\paragraph{表 3.1 统计物理学、贝叶斯推断、信息论及其认知解释}
最小化变分自由能，对应于最大化模型证据（或边缘似然）、最小化惊讶（或自信息）；其认知解释是知觉与行动。

最小化期望自由能，以及满足哈密顿最小作用量原理，对应于推断最可能的（或最不令人惊讶的）行动进程；其认知解释是“把规划理解为推断”。

达到非平衡稳态，对应于执行近似贝叶斯推断；其认知解释是自证。

能量函数上的梯度流，对应于模型证据上的梯度上升、惊讶上的梯度下降；其认知解释是神经元动力学。

\paragraph{框 3.2 统计物理学中的自由能与主动推断中的自由能}
自由能这一概念在统计物理学中被广泛使用，例如用于刻画热力学系统。虽然主动推断使用的是完全相同的方程，但它把这些方程应用于刻画主体的信念状态（相对于一个生成模型而言），而不是刻画其身体中的粒子之类的东西。因此，当我们说一个主动推断主体在最小化它的（变分）自由能时，我们指的是那些改变其信念状态的过程。为了避免误解，我们使用“变分自由能”这一术语，从而采用一种在机器学习中更常见的表述。

另一个更微妙的点在于，自由能的概念常常被用于平衡态统计热力学语境。主动推断所针对的却是生物体，也就是处于非平衡稳态、并且是开放的系统；它们与环境之间存在持续的、双向的交换。这是一个令人兴奋的新领域（Friston 2019a）。

\subsection{变分自由能、模型证据与惊讶}

第一个重要等价关系，是贝叶斯推断中的模型证据最大化（或边缘似然最大化）与变分自由能最小化之间的等价，而二者都在最小化惊讶。当我们诉诸一个特定的、用于处理推断难题的近似解法，即变分推断时，这种等价性就变得尤为明显。变分推断把推断问题改写为一个通过最小化自由能来求解的优化问题。自由能的最小值，对应于精确解近似得最好的那个点。把这种关系形式化写出来，就能看清这三个量之间的关系：

\begin{equation}
\Im(y\mid m) = -\ln P(y\mid m) \leq D_{KL}[Q(x)\,\|\,P(x\mid y,m)] - \ln P(y\mid m)
\tag{3.2}
\end{equation}

其中左边是惊讶，中间的 $-\ln P(y\mid m)$ 对应模型证据的负值，而右边整体对应变分自由能。

在式（3.2）中，与第 2 章不同，我们显式地把所有量都条件化在模型 $m$ 上，以强调这些量依赖于我们拥有的模型，或者说我们自身作为一个关于 $y$ 如何生成的模型；如果采用不同模型，这些量也会不同。这些量之间的等价性会引出一个问题：既然它们等价，为什么还要加以区分？主要原因在于，与模型证据不同，变分自由能可以被高效地最小化。回忆第 2 章，只有当 KL 散度项变为零时，变分自由能才与负模型证据或惊讶完全等价。这个条件并不总能严格达到，但可以逼近到非常接近零。因此，在为 $Q(x)$ 找到越来越好的取值的过程中，变分自由能也会越来越精确地逼近惊讶。我们已经多次提到这一点，因为必须反复强调：自由能与惊讶之间的核心关系，是本书的基础。具体地说，自由能是惊讶的一个上界。它可以与惊讶相等，也可以大于惊讶；大多少，由 KL 散度来量化。

这一点有一个有趣的方面：任何最小化其惊讶的系统，包括图 3.2 中那个非常简单的系统，也都在最小化某种自由能，只不过此时 $Q(x)$ 总是被设定为与精确后验概率相同，也就是把 KL 散度设为零。对于认知系统与非认知系统之间差异的一种看法是：后者总是具有零 KL 散度，而认知系统则必须先经过一个最小化这一项的（知觉）过程，然后它们的行动才会被保证能够最小化惊讶。注意，最小化这一散度，是知觉唯一能做的事情。这使内部状态的运动变得极其关键，因为这些内部状态所参数化的分布（图 3.2）必须尽可能接近精确后验。

然而，知觉无法最小化变分自由能中的第二个组成部分，也就是对应于实际惊讶的那个“证据”成分，因为它无法改变已经获得的感觉。只有通过行动，以改变感觉，主体才可能最小化变分自由能中的第二个成分，解决它的惊讶，或者等价地，最大化它的模型证据。这使得在自证过程中，给定内部状态时主动状态的运动变得至关重要。

一个例子有助于说明这一点。想象你的生成模型会根据饥饿水平预测血糖水平的分布，其中较高血糖对应饱腹，较低血糖对应饥饿。此外，假设这个模型赋予饱腹状态，也就是相对较高的血糖水平，以更高的先验概率，因此低血糖就变得令人惊讶。再设想你起初对自己的饥饿水平并不确定，却感知到低血糖。知觉会导向这样的推断：你是饥饿的，于是你体验到饥饿，这就闭合了 KL 散度。然而，知觉无法再进一步减少你的惊讶，无法缩小你先验上所期待的高血糖与实际感知到的低血糖之间的差距，因为知觉不能直接作用于你的感觉（低血糖）或其原因（生理过程）。你只有通过行动，去改变你所获取感觉的那个隐藏来源，才能最小化你的惊讶，例如吃一份甜点。

总而言之，知觉可以通过减少近似后验与真实后验之间的差异来最小化变分自由能，但它不能进一步最小化惊讶。惊讶最小化的下一步，是通过行动改变自己所获取的感觉；此时，推断超越了知觉，变得主动。

\subsection{期望自由能与最可能轨迹的推断}

另一个重要的等价关系，是期望自由能最小化与推断最可能的行动进程或策略之间的等价。这超出了仅仅规定状态空间中哪一部分最不令人惊讶的范围，而涉及抵达该区域或位置的不同路线本身可能有多令人惊讶。这些替代路径以策略的形式表达，而策略本质上就是跨越状态的轨迹。重要的是，在主动推断中，一项策略的对数概率被设定为与执行该策略时的期望自由能成比例。这意味着，最可能、或最不令人惊讶的路径，就是那个最小化期望自由能的路径。这一表述等价于物理学中对“作用量”的定义：作用量通过对某个能量沿路径进行积分（或求和），来给路径的概率打分。物理系统也许会面对一个假设路径空间，但它实际遵循的路径，是作用量最小的那一条，也就是哈密顿最小作用量原理。下一节将进一步展开这种主动推断与哈密顿最小作用量原理之间的类比。

\section{主动推断：理解行为与认知的一种新基础}

在最优控制、强化学习和经济学等领域，行为的优化来自状态的价值函数，并遵循贝尔曼方程（Sutton and Barto 1998）。本质上，每个状态（或状态--行动对）都会被赋予一个价值，这个价值代表该状态对主体来说有多“好”。状态（或状态--行动对）的价值通常通过试错学习得到，方法是统计从这些状态出发，在多少次尝试、以及经过多长时间之后，能够获得奖励。行为于是表现为通过到达高价值状态来优化奖励获取，也就是利用既往学习历史。

相比之下，在主动推断中，行为是推断的结果，而其优化是信念的函数。这一表述把（先验）信念与偏好结合在了一起。正如上文所述，采用期望自由能这一概念，相当于赋予主体一个隐含的先验信念：它将实现自己的偏好。因此，主体对于某条行动路线的偏好，就变成了它关于自己未来会做什么、未来会遇到什么的一种信念，或者说，关于自己将访问哪些未来状态轨迹的一种信念。这用（先验）信念的概念替代了价值的概念。

如果一个人具有强化学习背景，那么这种转换似乎相当奇怪，因为在强化学习中，价值与信念是分开的；如果一个人具有贝叶斯统计背景，这也会显得奇怪，因为在贝叶斯统计中，信念并不蕴含任何价值。然而，这是一种很有力量的转换，至少有三个原因。

第一，它自动包含了一个关于目的性行为（或目的论行为）的自洽过程模型，这与控制论的表述相近。如果我们赋予一个主动推断主体某种先验偏好，那么它就会行动去实现这种偏好，因为只有这样做，才与它关于“自己将会行动以实现其预期”的先验信念保持一致。注意，由此产生的（偏好）行动路线，也就是策略，在实验环境中是可以直接测量的；而价值函数或先验信念则需要被反推，因此是更间接、甚至近乎同义反复的测量。

第二，把行为看成信念（概率分布）的函数，会自动带来信念程度与不确定性之类的概念。这些概念支撑着适应性行动的重要方面，但在贝尔曼表述中并不能直接获得。出于同样的原因，这种表述在建模序列动力学和游走式行为时也更加灵活，而这些内容若仅用状态价值函数来建模会更困难（Friston, Daunizeau, and Kiebel 2009）。

第三，在这种表述下，最优行为会遵循统计物理学中的哈密顿最小作用量原理。事实上，主动推断在“行为是信念的函数”这一思想上又进一步：它还假定这种函数会成为一个能量函数，而主动推断主体最可能采取的行动路线，就是最小化自由能的那一条。一个深刻的结果是，生物体按照哈密顿最小作用量原理行事：它们沿着阻力最小的路径前进，直到到达稳态（或某条状态轨迹），图 3.3 中随机动力系统的行为就是例子。这是一个基础性假设，使主动推断不同于那些建立在贝尔曼表述之上的其他行为与认知理论。

这里值得简要说明一下，我们所说的哈密顿物理学与主动推断之间的类比是什么意思。这种类比包含三个层次。

第一，主动推断之于行为科学和生命科学的推进，可与拉格朗日表述和哈密顿表述之于牛顿力学的推进相比。牛顿力学最初是用微分方程来表述的，其中包括牛顿著名的第三定律，用来表达加速度与力之间的比例关系。后来，一种互补的视角出现了：人们开始通过考察动力系统中哪些量被守恒，来理解力学。于是，牛顿动力学也可以从这些守恒律中导出。这些守恒律提供了进一步理论推进的基础，并构成了随机物理学、相对论物理学和量子物理学某些部分的根基。类似地，主动推断重新表述了神经元动力学和行为动力学：过去这些动力学可能需要由一系列微分方程逐步搭建，而现在则可以通过指定一个量，也就是自由能，来导出这些动力学。正如不同种类的哈密顿量会导向不同类型的物理学一样，建立在不同生成模型基础上的自由能，也会导向不同的神经元动力学与行为动力学。

第二，哈密顿物理学与主动推断之间还有一种更直接的联系，即哈密顿量与概率测度之间的联系。这里的思想是把守恒的哈密顿量与系统的能量联系起来。请记住，到目前为止我们在这里以及第 2 章中所说的各种“能量”，都具有负对数概率的形式。这反映出一种把能量理解为“系统任何给定构型之不大可能性”的看法。从这个视角看，能量守恒与概率守恒是等价的定律。当耗散系统通过马尔可夫毯与外部状态耦合，并朝向低能量或高概率状态移动时，我们就可以直接把能量或哈密顿量与惊讶联系起来。因此，主动推断可以被看成应用于某一类特殊系统的哈密顿物理学，这类系统的特征就是拥有马尔可夫毯。

第三，这两种表述之间的联系还体现在把能量与动力学联系起来的变分法上。当哈密顿物理学被表达为最小作用量原理时，这一点尤其明显。这里的“作用量”指的是某个拉格朗日量沿一条路径的积分。关键在于，作用量是路径的泛函。这里，路径是一个时间函数，它在每个时刻给出该路径上粒子的位置与速度。一个确定性粒子实际遵循的路径，就是使该作用量最小的路径。类似地，主动推断建立在这样一个思想之上：信念，作为隐藏状态的函数，本身也必须最小化某个自由能泛函。这里真正的接触点在于：在两种情形中，被优化的都是函数（路径或信念），而相对于它们进行优化的都是泛函（分别是作用量或自由能）。因此，二者都属于变分法的范围，而变分法正是研究如何寻找泛函极值的数学分支。在物理学中，这导向欧拉--拉格朗日方程；在主动推断中，这导向变分推断程序。

\section{模型、策略与轨迹}

在第 3.2 节中，我们强调主动推断所适用的范围，是那些与环境之间享有某种分离的系统，并指出这在形式上体现为马尔可夫毯的存在。在第 3.3 节中，我们强调这种毯的持续存在要求系统具有这样一种动力学：平均而言最小化（感觉）状态的惊讶。由于这一点可以被解释为自证，我们就得到一个结论：行为由某种稳态分布所决定，而这种稳态分布又可以被解释为一个关于（感觉）数据如何生成的生成模型。

这告诉我们一件非常重要的事：每一种生成模型都应当对应不同类型的行为。因此，不同类型的行为可以通过指定不同生成模型来加以说明，也就是隐含地指定该系统会把什么视为令人惊讶的东西。进一步地，不同种类的生成模型还可能对应于具有不同复杂程度的适应性生物或认知生物（Corcoran et al. 2020）。像图 3.3 中驱动动力学的那类非常简单的生成模型，只提供了一种最小形式的认知，因为它们无法设想替代性的，或反事实性的，轨迹。此外，这些模型是浅层的，因为它们只能在单一时间尺度上进行推断。相比之下，层级式生成模型则允许在多个时间尺度上进行推断。在层级式或深层模型中，更高层级上的动力学通常编码那些变化较慢的事物，例如我正在读的这句话，而这些事物又为变化较快的事物，例如我正在读的这个词，提供语境，后者则由较低层级来表征（Kiebel et al. 2008; Friston, Parr, and de Vries 2017）。

那么，我们需要在一个模型中加入什么，才能导出那种更复杂、并且通常会与能动性和有感知系统联系在一起的行为呢？一个答案是：要具备对不同未来进行建模的能力，也就是对事件可能如何展开的不同方式进行建模，并在它们之间作出选择。反过来，考虑可能的未来，要求生成模型具有某种时间深度，并且显式表征行动的后果。把这些内容纳入模型，将会确保行为符合这些未来之中最可能实现的那个。具备这种对替代未来进行反事实思考的能力，或许正是把有感知系统所对应的稳态，与更简单生物区分开来的关键。当这些替代未来涉及我们能够控制的事物时，我们就把它们称为策略或计划。正如第 2 章所见，区分这些计划的一种方法，是把一种先验信念纳入模型：即那些具有最低期望自由能的策略，是最可信的策略。这为刻画某种带有马尔可夫毯、且处于稳态中的系统提供了一种方式，而这种刻画似乎与像我们这样的系统非常契合。

\section{在主动推断下对生成论、控制论与预测理论的统一}

通过强调自由能最小化，主动推断把三种表面上彼此脱节的理论视角统一并扩展起来。

第一，主动推断与关于生命和认知的生成论（enactive）理论是一致的。后者强调行为的自组织，以及与环境之间的自创生式交互，而这种交互确保生物体保持在可接受边界之内（Maturana and Varela 1980）。主动推断提供了一个形式框架，用来解释生物体如何通过自组织出一种统计结构，即马尔可夫毯，来抵抗其状态的离散化。这种结构既允许生物体与环境之间发生双向交换，同时又把生物体的状态与外部环境动力学分离开来，并在某种意义上保护了生物体状态的完整性。

第二，主动推断与控制论理论是一致的。控制论把行为描述为有目的的、目的论的。所谓目的论，意味着行为受到某种内部调节机制的调控，这种机制会持续检测目标是否达成；如果没有达成，就会引导纠正性行动（Rosenblueth et al. 1943, Wiener 1948, Ashby 1952, G. Miller et al. 1960, Powers 1973）。类似地，主动推断主体通过知觉与行动共同最小化偏好状态与感知状态之间的差异。主动推断通过明确指出真正被最小化的是一个主体可以测量的统计量，即变分自由能，来为这一最小化过程提供一种规范性的、可行的描述。在某些条件下，这个量对应于预测误差，也就是预期到的东西与实际感知到的东西之间的差别。这意味着，控制论控制可以被表述为一种前瞻性的过程，这也把我们引向下一点。

第三，主动推断与那些把控制描述为前瞻性过程的理论也是一致的，而这种前瞻性过程依赖于一个关于环境的模型，这个模型甚至可能在大脑中有物理实现（Craik 1943）。主动推断假定，主体使用一个（生成）模型来构造指导知觉与行动的预测，并评估其未来的、以及反事实的行动可能性。这一假定与“良好调节器定理”一致（Conant and Ashby 1970），该定理指出，任何控制器都应该拥有，或者本身就是，一个关于环境的良好模型。主动推断在严格的（近似）贝叶斯推断与（变分和期望）自由能最小化框架下，统一了这些基于模型的大脑与行为观点。进一步地，主动推断与观念运动理论（ideomotor theory）在很大程度上也是一致的（Herbart 1825, James 1890, Hoffmann 1993, Hommel et al. 2001）。观念运动理论认为，行动起始于一个想象过程，并且真正触发行动的，是对行动后果的一种预测性表征，而不是像刺激--反应理论那样由刺激来触发（Skinner 1938）。主动推断把这一思想置于一个推断框架中，在该框架里，行动源于一种关于未来的信念；这带来许多含义，例如，为了触发行动，主体必须暂时削弱感觉证据，因为否则这些证据会否证那种触发行动的信念（H. Brown et al. 2013）。

这些框架之间的统一很有意思，因为它们通常被认为彼此冲突。例如，在生物学中，自组织与目的论常常被视为不相容。此外，生成论理论往往弱化表征与控制的地位，而在大多数基于模型推断的理论中，表征与控制却是核心构件。主动推断从一个不寻常的角度对适应性主体的自创生动力学作出形式化，这个角度同时考虑了自组织与预测。通过连接不同视角，主动推断有可能帮助我们理解这些视角如何彼此照亮。

\section{主动推断：从生命的涌现到能动性}

主动推断从第一性原理出发，并逐步展开，用以解释从最简单到最复杂的适应系统和生命系统所表现出的行为与认知。在更简单生物与更复杂生物所构成的连续谱上，主动推断在两类系统之间划出一条界线：一类系统最小化变分自由能，另一类系统则还会最小化期望自由能。

任何通过主动采样感觉来最小化变分自由能的适应系统，等价地说，都是一个主动为其生成模型收集证据的主体，也就是一个自证主体（Hohwy 2016）。这类系统能够通过达成基本稳态过程所提供的设定点，来避免耗散、进行自我调节并存活下来。这些系统可以生成复杂而多样的行为形式，也可以具有很高的适应度，这一点在病毒身上已经可以看得很清楚。它们中的一些还可以拥有层级式生成模型，从而允许对在不同时间尺度上变化的事件进行推断，从更快的变化（位于较低层级）到更慢的变化（位于较高层级），并因此能够发展出复杂的策略来应对它们所经历的情境。然而，这些生物仍然有一个根本限制：它们的生成模型缺乏时间深度，也就是缺乏显式规划和显式考虑未来的能力，尽管它们也可以通过隐式方式做到这一点，例如作为遗传进化的结果。因此，它们始终活在当下。

一个具有时间深度的生成模型，为期望自由能最小化打开了大门，或者用心理学语言说，就是为规划打开了大门。在主动推断中，这带来的远不止是更强的适应性：它至少意味着一种原初形式的能动性。对于一个适应系统而言，最小化期望自由能，等价于拥有这样一种隐含先验：自己是一个最小化自由能的主体，而且会在未来通过行动来最小化自由能。当这种（先验）信念进入生成模型之后，适应系统就能够形成关于自己未来应当如何行动、将会追随哪些轨迹的信念。换句话说，它开始能够在不同的未来之间作出选择，而不只是像前文所述那些最简单的主体那样，仅仅选择如何应对已感知到的现在。因此，这种时间深度也转化为一种心理深度。至于不同生物如何分布在最简单与最复杂适应系统之间的连续谱上，以及它们能表现出什么形式的主动推断，这是一个经验问题。

\section{总结}

本章的主要内容可以概括如下：生物体必须确保自己只访问那些特征性的、或偏好的状态。如果把这些偏好状态定义为期望状态，那么就可以说，生物体必须最小化其感觉观测的惊讶，并维持一种最优熵，见框 3.3。

做到这一点，要求主体对环境动力学拥有某种自主性，并且配备一个马尔可夫毯，把它们的内部状态与环境的外部状态分开，也就是在二者之间表达一种条件独立。处于马尔可夫毯之内的主体，能够与环境进行双向的行动--知觉交换。这些交换由主动推断理论来形式化描述，在其中，知觉与行动都会最小化惊讶。之所以能够做到这一点，是因为主体配备了一个概率生成模型，用来说明其感觉观测如何生成。这个模型定义了惊讶，或者更准确地说，定义了一个可处理的代理量，即变分自由能；后者可以被测量，也可以被高效地最小化。

\paragraph{框 3.3 熵最小化与开放式行为}
主动推断建立在这样一个前提上：生物体努力维持相对的秩序（或负熵）、可控性和可预测性，尽管它们沉浸在一个其自然力量会不断制造波动、并带来无休止熵侵蚀威胁的环境中。这种主动追求秩序最基本的表现，就是生理稳态，因为若干关键的生理参数必须被维持在可生存区域之内。然而，熵最小化不应被等同于一套僵硬的反应库，例如自主神经式的稳态反应；恰恰相反，尤其在高级生物中，情况往往相反。我们可以发展出开放式的新行为库，以追求我们原初的稳态要求，例如为了满足口渴和其他需要而生产并购买好酒。这有时被称为“异稳态”（allostasis）（Sterling 2012）。

更广泛地说，我们还会主动追求某种秩序与可控性本身，而不一定直接指向某个具体的稳态要求，或许是因为保留秩序本身就有助于满足许多这类要求。我们会主动雕刻自己的生态位，使之变得更可预测、更不令人惊讶。这一点体现在我们构造物理空间的方式中，例如庇护所和城市，它们使我们免受失控自然力量的侵袭；也体现在我们构造文化空间的方式中，例如拥有法律和义务规范的社会，它们使我们免受无政府状态社会力量的侵袭。在所有这些例子里，我们通常都需要接受某种短期的熵或惊讶增加，例如建造新事物或改变社会立场，以确保长期的熵或惊讶下降。这帮助我们理解，惊讶最小化这一基本要求并不与认知性要求、新奇寻求、好奇和探索行为相冲突；相反，它恰恰促进了这些被我们视为许多物种核心特征的行为。

认知性要求首先体现在变分自由能最小化过程中。对自由能的一种分解方式，是把它表达为在近似后验下取期望的 Gibbs 能量减去近似后验的熵。换言之，主体是在努力增加熵。虽然这看起来像个悖论，但一旦注意到这里增加的是主体信念，也就是其近似后验信念的熵，这个悖论就消失了。可以把它理解为这样一种要求：既要尽可能准确地解释事物，又要“保持选项开放”，除非必要，不要过早承诺于任何一个具体解释，这就是最大熵原理（Jaynes 1957）。

认知性动力学表现出来的第二种方式，是在期望自由能最小化过程中。有意思的是，在这里会出现两个符号相反的熵。它们包括：后验预测熵，也就是“在给定某种选择时，我对自己会遇到什么结果有多不确定”，这一项必须被最大化，这一点与变分自由能中关于状态信念的情况相同；以及在给定状态条件下结果的条件熵，也就是某项策略所包含的歧义性，这一项必须被最小化。因此，在变分自由能最小化过程中，要求是最大化关于当前信念的熵；而在期望自由能最大化过程中，要求则是选择那些能够最小化未来信念歧义性的行动。由此便产生了认知性的、好奇的、寻求新奇的、搜寻信息的行为，它们支持不确定性的消解，或者支持生成模型的改进，而这反过来又会在长期中最小化惊讶（Seth 2013; Friston, Rigoli et al. 2015; Seth and Friston 2016; Schwartenbeck, Passecker et al. 2019）。

一个主动推断主体看起来是在某个生成模型之下执行（近似）贝叶斯推断，并为自己的模型最大化证据，也就是说，它是一个自证主体。推断中的前瞻性部分，则通过选择那些预期将在未来最小化自由能的行动路线或策略来实现。这一形式体系导向了一种关于（最优）行为的新视角：即把行为理解为哈密顿最小作用量原理的体现。这是一条把主动推断与统计物理学、热力学以及非平衡稳态领域连接起来的第一性原理。
